\documentclass[11pt]{article}

\usepackage[margin=1.1in]{geometry}
\usepackage{amsmath,amssymb,amsthm,amscd}
\usepackage{booktabs}
\usepackage{microtype}
\usepackage{xcolor}
\usepackage[colorlinks=true,linkcolor=blue!55!black,citecolor=blue!55!black,urlcolor=blue!55!black]{hyperref}

% ---------- macros ----------
\newcommand{\cT}{\mathcal{T}}   % planar binary trees (free magma)
\newcommand{\cW}{\mathcal{W}}   % words (free semigroup)
\newcommand{\cB}{\mathcal{B}}   % unordered binary trees (free commutative magma)
\newcommand{\cM}{\mathcal{M}}   % prime multisets (free commutative semigroup)
\newcommand{\cG}{\mathcal{G}}   % grouped quantities
\newcommand{\cA}{\mathcal{A}}   % flat quantities
\newcommand{\NN}{\mathbb{N}}
\newcommand{\pkg}[1]{\boxed{#1}}
\newcommand{\act}{\rho}
\DeclareMathOperator{\supp}{supp}
\numberwithin{equation}{section}

\theoremstyle{plain}
\newtheorem{theorem}{Theorem}[section]
\newtheorem{proposition}[theorem]{Proposition}
\newtheorem{lemma}[theorem]{Lemma}
\newtheorem{corollary}[theorem]{Corollary}
\newtheorem{conjecture}[theorem]{Conjecture}
\theoremstyle{definition}
\newtheorem{definition}[theorem]{Definition}
\newtheorem{example}[theorem]{Example}
\newtheorem{question}[theorem]{Question}
\theoremstyle{remark}
\newtheorem{remark}[theorem]{Remark}

\title{Prime-Leaf Tree Theory\\[2mm]
\large Ordinary arithmetic as the shadow of a free multiplicative structure}
\author{Giuseppe Barbalinardo}
\date{Working draft --- \today}

\begin{document}
\maketitle

\begin{abstract}
We develop an arithmetic in which multiplication does not return a number.
The product of two objects is a \emph{multiplication tree}: an ordered, parenthesized
record of how the object was built from \emph{prime leaves} $L_2, L_3, L_5, \dots$
Formally, the multiplicative world is the free magma $\cT$ on the set of primes, and
the ordinary value of a tree is recovered by a shadow map $N \colon \cT \to \NN_{\ge 2}$,
the unique magma morphism sending $L_p$ to $p$. Ordinary multiplicative arithmetic is
the double quotient of $\cT$ by associativity and commutativity; the Fundamental
Theorem of Arithmetic is precisely the statement that this double quotient loses no
information. The two identifications are independent, giving a commuting square of four
arithmetics, each with a degeneracy count and a zeta function: the Riemann zeta function
(multisets), a geometric zeta with critical exponent $1.39943\dots$ (words), an
Otter-type zeta with exponent $1.98847\dots$ (unordered trees), and a Catalan zeta with
exponent $\sigma^{*} = 2.59738\dots$ (planar trees), where $P(\sigma^{*}) = \tfrac14$ for
the prime zeta function $P$. The tree partition function converges \emph{at} its critical
point with the exact value $\tfrac12$ and a square-root cusp: a second-order phase
transition, in contrast with the first-order pole of $\zeta$. Its complex branch points
--- the ``tree zeros'', solutions of $P(s) = \tfrac14$ --- are computed; they do not lie
on a vertical line. Finally we construct the bridge between the multiplicative and
additive worlds: trees do not map into quantities, they \emph{act} on them, producing
grouped quantities. Under this action, associativity of multiplication becomes
additively invisible, commutativity survives grouping and dies only under flattening,
and distributivity becomes a theorem about flattening rather than an axiom. The prime
leaves themselves can then be derived, rather than imported: primes are exactly the heap
sizes that admit no uniform regrouping.
\end{abstract}

\tableofcontents

%======================================================================
\section{Introduction}

\subsection{Multiplication before evaluation}

Ordinary arithmetic writes
\[
3 \times 2 = 6
\qquad\text{and}\qquad
1+1+1+1+1+1 = 6
\]
and treats the two results as the same kind of object. The starting intuition of this
paper is that they are not. Six liters obtained by pouring is a flat additive quantity.
``Three groups of two'' is a \emph{structured} object --- and a differently structured
object than ``two groups of three'' --- whose ordinary numerical shadow happens to be~$6$.
The slogan:

\begin{quote}
\emph{You never get six liters from multiplication. You get a structured object whose
shadow is $6$.}
\end{quote}

We therefore refuse to let multiplication collapse into numbers. Multiplication forms
\emph{trees}: the product of two trees is the ordered pair of them, written
$T_1 \boxtimes T_2$, and nothing is evaluated. A tree is an unevaluated multiplication
expression; the familiar integer is obtained only when one applies an evaluator
$N$, the \emph{shadow map}. In this reading the classical laws --- commutativity,
associativity, distributivity --- are not properties of multiplication itself. They are
properties of the \emph{evaluator}: statements that $N$ merges distinct expressions.
The richer theory keeps the expressions apart and studies what the merging forgets.

\subsection{Results}

The paper makes this program precise and proves the following.

\begin{enumerate}
\item \textbf{The framework} (\S\ref{sec:trees}). The multiplicative world is the free
magma $\cT$ on the set of primes; the shadow map is the unique magma morphism
$N\colon \cT \to (\NN_{\ge2},\cdot)$ with $N(L_p) = p$. The multiplicative world has no
unit: $1$ is not the shadow of any tree, and appears only after quotienting, as the
empty multiset.

\item \textbf{The quotient square} (\S\ref{sec:square}). Associativity and commutativity
are \emph{independent} identifications. Quotienting $\cT$ by either, both, or neither
yields four arithmetics
\[
\cT \;(\text{planar trees}), \quad
\cW \;(\text{words}), \quad
\cB \;(\text{unordered trees}), \quad
\cM \;(\text{multisets}),
\]
arranged in a commuting square. The Fundamental Theorem of Arithmetic is the statement
that the shadow map of the final quotient $\cM$ is a bijection onto $\NN_{\ge2}$. Each
level carries a degeneracy count over each integer $n$; for $\cT$ it is
$C_{\Omega(n)-1}\,\Omega(n)!/\prod_p \alpha_p!$ (Catalan number times multinomial), for
$\cW$ the multinomial, for $\cB$ a sequence which restricts to the
Wedderburn--Etherington numbers on prime powers and which appears not to have been
recorded before, and for $\cM$ the constant $1$. All counts were verified by exhaustive
enumeration for $n \le 20000$ (Appendix~\ref{app:verification}).

\item \textbf{Four zeta functions} (\S\ref{sec:zeta}). Each corner has a partition
function $Z(s) = \sum N(\,\cdot\,)^{-s}$. They satisfy, respectively, a Catalan equation
$Z = P + Z^2$, a geometric equation $Z = P/(1-P)$, an Otter equation
$Z = P + \tfrac12(Z^2 + Z(2s))$, and the Euler product
$\zeta(s) = \exp\bigl(\sum_{k\ge1} P(ks)/k\bigr)$, where
$P(s) = \sum_p p^{-s}$ is the prime zeta function. The critical exponents form a strict
ladder
\[
1 \;<\; 1.3994333287263303 \;<\; 1.9884768518009081 \;<\; \sigma^{*} = 2.5973851271346716,
\]
one for each amount of remembered structure. At its critical point the tree zeta
converges with the exact value $Z_\cT(\sigma^{*}) = \tfrac12$ and a square-root cusp.

\item \textbf{Tree zeros} (\S\ref{sec:analytic}). $Z_\cT$ is an algebraic function of
$P$, so its analytic structure reduces to that of $P$ --- which inherits a branch point
at $s = \rho/k$ from every nontrivial zero $\rho$ of $\zeta$ and has the imaginary axis
as a natural boundary. The genuinely new singularities are the branch points where
$P(s) = \tfrac14$. We compute the five with $0 < \operatorname{Im} s < 50$; they cluster
near integer multiples of $2\pi/\log 2$ and do \emph{not} lie on a vertical line, so the
naive tree analogue of the Riemann Hypothesis is false and must be reformulated as a
distribution question (\S\ref{sec:open}).

\item \textbf{Thermodynamics} (\S\ref{sec:thermo}). With energy $E(T) = \log N(T)$, the
four zetas are partition functions of four gases built on the primon gas of Julia and of
Bost--Connes. Remembering more structure raises the critical inverse temperature ---
equivalently, lowers the Hagedorn temperature from $1$ to $0.38500\dots$ --- and changes
the order of the transition: the associative levels blow up at their critical point
(pole), the non-associative levels stay finite and cusp (continuous transition).

\item \textbf{The additive bridge} (\S\ref{sec:bridge}). Multiplication trees do not map
into the additive world; they \emph{act} on it. We build a world $\cG$ of grouped
quantities (boxed heaps of unit quantities), on which a tree $T$ acts by uniform
regrouping. Three theorems then stratify the classical laws:
associativity of $\boxtimes$ is invisible to the action (the action factors through the
word quotient, \emph{exactly}); commutativity fails for the action and holds only after
flattening; distributivity fails structurally and holds after flattening. ``You never
get six liters from multiplication'' becomes a theorem. Moreover the prime leaves can be
\emph{derived} from the additive world: primes are exactly the heap sizes admitting no
uniform regrouping, which answers the foundational question of whether the leaf alphabet
must be imported.
\end{enumerate}

\subsection{Relation to existing mathematics}\label{sec:priorart}

Parts of this structure are classical, and it is worth saying precisely which parts.
The free magma and its Catalan combinatorics are standard. The count of ordered
factorizations into primes is OEIS A008480, and the full planar-tree degeneracy
$g_\cT(n)$ appears as OEIS A144757 (``number of factor trees for $n$'')~\cite{oeis}. The
thermodynamic reading of $\zeta$ as a partition function over multiplicative states is
the primon gas of Julia~\cite{julia} and Spector~\cite{spector}, made deep by
Bost--Connes~\cite{bostconnes}. Closest in spirit is Loday's
\emph{arithmetree}~\cite{loday}, further developed by Bruno and
Yasaki~\cite{brunoyasaki}: an arithmetic on planar binary trees into which $\NN$ embeds,
with non-commutative addition and one-sided distributivity. The ordered-factorization
growth problem with arbitrary parts goes back to Kalm\'ar~\cite{kalmar}, and the natural
boundary of the prime zeta function is due to Landau and Walfisz~\cite{landauwalfisz}.

Against this background, the contributions here are: the quotient \emph{square} (rather
than a single chain) with the new unordered-tree corner, its Otter zeta and critical
exponent; the exact critical behavior $Z_\cT(\sigma^{*}) = \tfrac12$ and the transition-order
dichotomy; the computation and non-collinearity of the tree zeros; the grouped-quantity
bridge with the stratification theorems; and the internal derivation of the prime
leaves. One honest caution runs through the paper: a \emph{free} structure carries no
secrets about the primes themselves. Every classically hard question enters either
through the values $2,3,5,\dots$ fed to the shadow map --- that is, through $P(s)$ ---
or through the additive bridge. What the theory offers is a clean separation of the
structural from the arithmetic, and new invariants that live in the separated layers.

%======================================================================
\section{The multiplicative world}\label{sec:trees}

\subsection{Prime leaves and multiplication trees}

\begin{definition}[Multiplication trees]\label{def:trees}
Fix the set of prime leaves $\{L_p : p \text{ prime}\}$. The set $\cT$ of
\emph{multiplication trees} is defined by the grammar
\[
T \;::=\; L_p \;\mid\; (T_1 \boxtimes T_2).
\]
Thus $\cT$ is the \emph{free magma} on the set of primes: $\boxtimes$ is ordered-pair
formation, subject to no laws whatsoever.
\end{definition}

Examples of distinct elements of $\cT$:
\[
L_2, \qquad
L_2 \boxtimes L_3, \qquad
L_3 \boxtimes L_2, \qquad
(L_2 \boxtimes L_3) \boxtimes L_5, \qquad
L_2 \boxtimes (L_3 \boxtimes L_5).
\]
Non-commutativity and non-associativity are not axioms we impose; they are the
\emph{absence} of axioms. $L_2 \boxtimes L_3$ records ``a structure built from $L_2$,
then $L_3$''; the parenthesization records construction history. A multiplication tree
is best read as the parse tree of a multiplication expression, before evaluation.

For $T \in \cT$ write $|T|$ for the number of leaves and $\mathrm{word}(T)$ for the
left-to-right sequence of leaf labels; $\supp(T)$ denotes the set of primes occurring.

\subsection{The shadow map}

\begin{definition}[Shadow]\label{def:shadow}
The \emph{shadow map} $N \colon \cT \to \NN_{\ge 2}$ is defined recursively by
\[
N(L_p) = p, \qquad N(T_1 \boxtimes T_2) = N(T_1)\,N(T_2).
\]
Equivalently: $N$ is the unique magma morphism from $\cT$ to the multiplicative
semigroup $(\NN_{\ge2},\cdot)$ sending $L_p$ to $p$. It is the evaluator of
multiplication expressions.
\end{definition}

\begin{remark}[No unit]\label{rem:nounit}
$\cT$ has no unit object and $1$ is not the shadow of any tree: the image of $N$ is
exactly $\NN_{\ge 2}$. One could freely adjoin an empty tree $I$, but then
$I \boxtimes T \neq T$ (freeness again), which buys nothing. We prefer to let $1$ appear
where it belongs: as the \emph{empty multiset} after quotienting
(\S\ref{sec:square}), i.e.\ as a feature of the collapsed theory, not of the fundamental
one.
\end{remark}

Composite numbers are not primitive objects of the multiplicative world. There is no
$L_6$; the number $6$ occurs only as the common shadow of the two trees
$L_2\boxtimes L_3$ and $L_3\boxtimes L_2$. Likewise $8$ is the common shadow of
$(L_2\boxtimes L_2)\boxtimes L_2$ and $L_2\boxtimes(L_2\boxtimes L_2)$. Distinct
microstates, one observable: this degeneracy is quantified in \S\ref{sec:square} and
made thermodynamic in \S\ref{sec:thermo}.

%======================================================================
\section{Ordinary arithmetic as a quotient: the square}\label{sec:square}

\subsection{Two independent identifications}

\begin{definition}\label{def:congruences}
Let $\approx_a$ be the smallest congruence on $(\cT,\boxtimes)$ with
$(T_1\boxtimes T_2)\boxtimes T_3 \approx_a T_1\boxtimes(T_2\boxtimes T_3)$ for all
$T_i$, and $\approx_c$ the smallest congruence with
$T_1\boxtimes T_2 \approx_c T_2\boxtimes T_1$. Write
\[
\cW = \cT/\!\approx_a, \qquad
\cB = \cT/\!\approx_c, \qquad
\cM = \cT/(\approx_a \vee \approx_c).
\]
\end{definition}

Each quotient is a familiar free object: $\cW$ is the free semigroup on the primes
(elements are nonempty \emph{words} in primes, i.e.\ ordered prime factorizations);
$\cB$ is the free \emph{commutative} magma on the primes (binary trees with unordered
branches); $\cM$ is the free commutative semigroup (nonempty finite \emph{multisets} of
primes). The shadow map descends to all four corners, and the two identifications
commute:
\[
\begin{CD}
\cT @>\text{associativity}>> \cW \\
@V\text{commutativity}VV @VV\text{commutativity}V \\
\cB @>>\text{associativity}> \cM
\end{CD}
\]
The shadow map of the final corner is where classical arithmetic begins:

\begin{proposition}[FTA as exactness of the collapse]\label{prop:fta}
The induced map $N \colon \cM \to \NN_{\ge2}$ is a bijection. This statement is
equivalent to the Fundamental Theorem of Arithmetic.
\end{proposition}

\begin{proof}
Elements of $\cM$ are nonempty prime multisets $\{p_1,\dots,p_k\}$ with
$N = p_1\cdots p_k$. Surjectivity of $N$ is the existence of a prime factorization;
injectivity is its uniqueness. \end{proof}

Thus classical multiplicative arithmetic is the coarsest shadow of the tree theory, and
the FTA acquires a new reading: \emph{the double quotient loses nothing numerically}.
Everything it does lose is structure, and the degeneracy counts below measure exactly
how much.

\begin{remark}[The identifications are themselves structured]
The associativity identification is not a group action: two trees with the same word are
connected by a sequence of reassociations, and the set of parenthesizations of a fixed
word carries the \emph{Tamari lattice} structure, realized geometrically by the
associahedron~\cite{tamari}. Commutativity moves similarly generate the permutohedron.
The quotient square therefore has polytopal fibers, a refinement we do not pursue here.
\end{remark}

\subsection{Degeneracy counts}

For $n \ge 2$ with factorization $n = \prod_p p^{\alpha_p}$ write
$\Omega(n) = \sum_p \alpha_p$. For each corner $\mathcal{X}$ of the square let
$g_{\mathcal X}(n) = \#\{X \in \mathcal{X} : N(X) = n\}$ be the number of microstates
over $n$.

\begin{proposition}[Planar trees]\label{prop:gT}
\[
g_\cT(n) \;=\; C_{\Omega(n)-1}\;\frac{\Omega(n)!}{\prod_p \alpha_p!},
\qquad C_k = \frac{1}{k+1}\binom{2k}{k}.
\]
\end{proposition}

\begin{proof}
A tree with shadow $n$ has exactly $\Omega(n)$ leaves, by induction and multiplicativity
of $N$. The map $T \mapsto (\text{shape}(T), \mathrm{word}(T))$ is a bijection from
$\{T : N(T)=n\}$ to (binary tree shapes with $k=\Omega(n)$ leaves) $\times$ (sequences
of primes with product $n$). There are $C_{k-1}$ shapes~\cite{stanley} and
$\Omega!/\prod \alpha_p!$ sequences (permutations of the prime multiset). The two
factors are independent: the fiber of the full collapse $\cT \to \cM$ splits as
shapes $\times$ orderings.
\end{proof}

\begin{proposition}[Words]\label{prop:gW}
$g_\cW(n) = \Omega(n)!/\prod_p \alpha_p!$, with recursion
$g_\cW(n) = \sum_{p \mid n} g_\cW(n/p)$ and $g_\cW(p)=1$. This is OEIS A008480.
\end{proposition}

\begin{proposition}[Unordered trees]\label{prop:gB}
$g_\cB$ satisfies $g_\cB(p) = 1$ and, for composite $n$,
\[
g_\cB(n) \;=\; \frac12 \sum_{\substack{d \mid n \\ 1<d<n}} g_\cB(d)\, g_\cB(n/d)
\;+\; \frac12\,[\,n \text{ is a square}\,]\; g_\cB(\sqrt n).
\]
On prime powers, $g_\cB(p^k)$ is the $k$-th Wedderburn--Etherington number
(OEIS A001190): 1, 1, 1, 2, 3, 6, 11, 23, \dots
\end{proposition}

\begin{proof}
An element of $\cB$ over composite $n$ is an unordered pair $\{B_1,B_2\}$ with
$N(B_1)N(B_2) = n$. Ordered pairs are counted by $\sum_{1<d<n} g_\cB(d)g_\cB(n/d)$; the
diagonal (possible only when $n$ is a square, both factors of shadow $\sqrt n$ and equal
as trees) is counted by $g_\cB(\sqrt n)$; unordered pairs with repetition are (ordered
$+$ diagonal)$/2$. On $p^k$ all leaves are identical and the count reduces to
unordered binary trees with $k$ indistinguishable leaves, the defining recursion of the
Wedderburn--Etherington numbers.
\end{proof}

\begin{remark}[A sequence not in the OEIS]\label{rem:newseq}
At the time of writing, the sequence $g_\cB(2), g_\cB(3), \dots =
1,1,1,1,1,1,1,1,1,1,2,\dots$ does not appear in the OEIS. The closest entry, A375120
(``complete binary unordered tree-factorizations''), uses a different convention: it
counts the two children as an \emph{ordered} pair whenever the factor split is
$(d,d)$, so it disagrees with Proposition~\ref{prop:gB} first at $n=144$ (it gives
$42$ where the true unordered count is $41$) and disagrees with
Wedderburn--Etherington on prime powers. Submitting $g_\cB$ is a small item of
independent interest.
\end{remark}

\begin{example}[The four levels over $n = 12$]\label{ex:twelve}
\[
\begin{array}{lll}
\cT: & 6 \text{ trees} &
(L_2\boxtimes L_2)\boxtimes L_3,\;
(L_2\boxtimes L_3)\boxtimes L_2,\;
(L_3\boxtimes L_2)\boxtimes L_2, \\
& & L_2\boxtimes(L_2\boxtimes L_3),\;
L_2\boxtimes(L_3\boxtimes L_2),\;
L_3\boxtimes(L_2\boxtimes L_2); \\[1mm]
\cW: & 3 \text{ words} & (2,2,3),\; (2,3,2),\; (3,2,2); \\[1mm]
\cB: & 2 \text{ bushes} & \{\{L_2,L_2\},L_3\},\; \{\{L_2,L_3\},L_2\}; \\[1mm]
\cM: & 1 \text{ multiset} & \{2,2,3\}.
\end{array}
\]
And $g_\cT(12) = C_2 \cdot \tfrac{3!}{2!\,1!} = 2\cdot 3 = 6$, as predicted.
\end{example}

Table~\ref{tab:degeneracies} lists the three nontrivial counts for small $n$. All
formulas and recursions of this subsection were verified against brute-force enumeration
for every $n \le 20000$ (Appendix~\ref{app:verification}).

\begin{table}[ht]
\centering\small
\begin{tabular}{@{}rrrrr@{\hspace{9mm}}rrrrr@{}}
\toprule
$n$ & $\Omega$ & $g_\cT$ & $g_\cB$ & $g_\cW$ &
$n$ & $\Omega$ & $g_\cT$ & $g_\cB$ & $g_\cW$ \\
\midrule
 2 & 1 &  1 & 1 & 1 &  20 & 3 &  6 & 2 & 3 \\
 4 & 2 &  1 & 1 & 1 &  24 & 4 & 20 & 4 & 4 \\
 6 & 2 &  2 & 1 & 2 &  27 & 3 &  2 & 1 & 1 \\
 8 & 3 &  2 & 1 & 1 &  28 & 3 &  6 & 2 & 3 \\
10 & 2 &  2 & 1 & 2 &  30 & 3 & 12 & 3 & 6 \\
12 & 3 &  6 & 2 & 3 &  32 & 5 & 14 & 3 & 1 \\
16 & 4 &  5 & 2 & 1 &  36 & 4 & 30 & 6 & 6 \\
18 & 3 &  6 & 2 & 3 &  64 & 6 & 42 & 6 & 1 \\
\bottomrule
\end{tabular}
\caption{Degeneracies over $n$ at the three structured levels (primes and prime-free
rows with all counts $1$ omitted). Note $g_\cT = C_{\Omega-1}\cdot g_\cW$ throughout,
and $g_\cB(2^k) = 1,1,1,2,3,6$ for $k=1,\dots,6$ (Wedderburn--Etherington).}
\label{tab:degeneracies}
\end{table}

%======================================================================
\section{Four zeta functions}\label{sec:zeta}

\subsection{Definitions and functional equations}

Let $P(s) = \sum_p p^{-s}$ denote the prime zeta function. For each corner
$\mathcal{X} \in \{\cT, \cW, \cB, \cM\}$ define the partition function
\[
Z_{\mathcal X}(s) \;=\; \sum_{X \in \mathcal X} \frac{1}{N(X)^{s}}
\;=\; \sum_{n \ge 2} \frac{g_{\mathcal X}(n)}{n^{s}}
\]
(for $\cM$ we add by hand the empty multiset, of shadow $1$, so that
$Z_\cM = \zeta$ exactly).

\begin{theorem}[Functional equations]\label{thm:funceq}
As identities of Dirichlet series,
\begin{align}
Z_\cT(s) &= P(s) + Z_\cT(s)^2,
& Z_\cT &= \frac{1 - \sqrt{1 - 4P(s)}}{2} = \sum_{k\ge1} C_{k-1}\,P(s)^k,
\label{eq:cat}\\
Z_\cW(s) &= P(s) + P(s)\,Z_\cW(s),
& Z_\cW &= \frac{P(s)}{1-P(s)} = \sum_{k \ge 1} P(s)^k,
\label{eq:geo}\\
Z_\cB(s) &= P(s) + \tfrac12\bigl(Z_\cB(s)^2 + Z_\cB(2s)\bigr),
& Z_\cB &= 1 - \sqrt{1 - 2P(s) - Z_\cB(2s)},
\label{eq:otter}\\
Z_\cM(s) &= \zeta(s) = \prod_p \frac{1}{1-p^{-s}}
& \zeta &= \exp\Bigl(\sum_{k \ge 1} \frac{P(ks)}{k}\Bigr).
\label{eq:euler}
\end{align}
In \eqref{eq:cat} and \eqref{eq:otter} the branch of the square root is fixed by
$Z \to 0$ as $\operatorname{Re} s \to \infty$.
\end{theorem}

\begin{proof}
\eqref{eq:cat}: a tree is a leaf or an ordered pair of trees, and $N^{-s}$ is
multiplicative over $\boxtimes$, so the sum over pairs factors as $Z_\cT^2$. Solving the
quadratic and expanding gives the Catalan series.
\eqref{eq:geo}: a word is a prime or a prime followed by a word.
\eqref{eq:otter}: an element of $\cB$ is a leaf or an unordered pair
$\{B_1,B_2\}$, possibly equal; ordered pairs contribute $Z_\cB^2$, the diagonal
contributes $\sum_B N(B)^{-2s} = Z_\cB(2s)$, and unordered pairs with repetition are
half their sum --- the same P\'olya computation as in Proposition~\ref{prop:gB}.
\eqref{eq:euler}: independence of leaf multiplicities gives the Euler product, and
$\log \zeta(s) = \sum_p \sum_{k\ge1} p^{-ks}/k = \sum_k P(ks)/k$.
\end{proof}

\begin{remark}[One zeta per way of assembling primes]\label{rem:species}
The four right-hand sides are the four classical constructions of the symbolic
method~\cite{flajolet}: sequence, multiset, planar binary tree, unordered binary tree,
each applied to the ``atomic'' Dirichlet element $P$. The answer to the question
\emph{which generalized zeta function is the right one} is therefore: there is one for
every combinatorial construction, all expressible in $P(s), P(2s), P(3s), \dots$, and
the Riemann zeta function is the multiset member of the family. Note where the unit
lives: only the multiset construction has an empty structure, which is why $n=1$ belongs
to $\zeta$ and to no other corner (Remark~\ref{rem:nounit}).
\end{remark}

\subsection{Critical exponents}

\begin{theorem}[The exponent ladder]\label{thm:ladder}
Each series has a finite abscissa of convergence $\sigma_{\mathcal X}$, determined by
$P$:
\[
\begin{array}{llll}
\cM: & \sigma_\cM = 1 & \text{(pole of } \zeta), & Z \to \infty; \\[0.5mm]
\cW: & P(\sigma_\cW) = 1, & \sigma_\cW = 1.3994333287263303\dots, & Z \to \infty; \\[0.5mm]
\cB: & 2P(\sigma_\cB) + Z_\cB(2\sigma_\cB) = 1, & \sigma_\cB = 1.9884768518009081\dots, & Z \to 1; \\[0.5mm]
\cT: & P(\sigma^{*}) = \tfrac14, & \sigma^{*} = 2.5973851271346716\dots, & Z \to \tfrac12,
\end{array}
\]
and $1 < \sigma_\cW < \sigma_\cB < \sigma^{*}$: every remembered layer of structure
strictly raises the critical exponent.
\end{theorem}

\begin{proof}[Proof sketch]
All coefficients are nonnegative, so each abscissa is located by the real singularity.
For $\cW$, the geometric series diverges exactly when $P \ge 1$. For $\cT$, the Catalan
series $\sum C_{k-1}P^k$ has radius of convergence $\tfrac14$ in $P$; since
$C_{k-1}4^{-k} \sim k^{-3/2}/(4\sqrt{\pi})$ is summable, the series still converges at
$P = \tfrac14$, with value $\tfrac{1-\sqrt{1-1}}{2}$ evaluated as the limit
$\sum_k C_{k-1}4^{-k} = \tfrac12$. For $\cB$, equation \eqref{eq:otter} defines $Z_\cB$
by descending the argument-doubling cascade ($Z_\cB(2s)$ lies deep in the convergent
region), and the singularity occurs where the discriminant $1 - 2P(s) - Z_\cB(2s)$
vanishes, at which point $Z_\cB = 1$: this is Otter's criticality condition for
unordered trees~\cite{otter,flajolet}. The numerical values are computed to the stated
precision in Appendix~\ref{app:verification}; the strict ordering follows from
$g_\cM \le g_\cW, g_\cB \le g_\cT$ together with strict inequality on a positive-density
set of $n$.
\end{proof}

\begin{remark}[Behavior at criticality: pole versus cusp]\label{rem:cusp}
The associative levels diverge at their critical point: $\zeta$ and $Z_\cW$ have simple
poles ($Z_\cW$ because $P'(\sigma_\cW) = -1.7311562\dots \ne 0$). The non-associative
levels converge at theirs: expanding \eqref{eq:cat} at $\sigma^{*}$,
\[
Z_\cT(s) \;=\; \tfrac12 - \sqrt{\,|P'(\sigma^{*})|\,}\,(s-\sigma^{*})^{1/2} + O(s-\sigma^{*}),
\qquad \sqrt{|P'(\sigma^{*})|} = 0.4796402876\dots,
\]
a square-root cusp with a finite value and an infinite derivative, and similarly for
$Z_\cB$ at $\sigma_\cB$ with $Z_\cB = 1$. Associativity produces geometric/exponential
(pole-type) generating structures; its absence produces Catalan/Otter (branch-type)
ones. The thermodynamic meaning of this dichotomy is taken up in \S\ref{sec:thermo}.
\end{remark}

\begin{remark}[Growth of the counting functions]\label{rem:growth}
By the Wiener--Ikehara theorem applied to $Z_\cW$ (legitimate: the line
$\operatorname{Re} s = \sigma_\cW$ lies in the half-plane of analyticity of $P$, and
$|P(\sigma_\cW + it)| < 1$ for $t \ne 0$ by Lemma~\ref{lem:strict} below),
\[
\sum_{n \le x} g_\cW(n) \;\sim\; \frac{x^{\sigma_\cW}}{\sigma_\cW\,|P'(\sigma_\cW)|}
\;\approx\; 0.4128\, x^{1.39943\dots}.
\]
For the tree level, transfer heuristics for the square-root singularity suggest
$\sum_{n \le x} g_\cT(n) \asymp x^{\sigma^{*}} (\log x)^{-3/2}$; making this rigorous is
Problem~\ref{prob:tauberian}. Note also the extreme values: along $n = 2^k$ one has
$g_\cT(n) = C_{k-1}$, which grows like $n^{2}/(\log n)^{3/2}$ up to constants --- the
microstate count over a single observable can approach the \emph{square} of the
observable. For comparison, Kalm\'ar's classical problem~\cite{kalmar} of ordered
factorizations with arbitrary parts $\ge 2$ has Dirichlet series $1/(2 - \zeta(s))$ and
growth exponent $\rho = 1.7286472389981836\dots$ with $\zeta(\rho) = 2$; it sits between
our word and bush levels but belongs to a different family (its parts are composite-friendly).
\end{remark}

%======================================================================
\section{Analytic structure: inherited singularities and tree zeros}\label{sec:analytic}

Since $Z_\cT = \tfrac12(1 - \sqrt{1-4P})$ is an \emph{algebraic} function of $P$, its
analytic life is entirely inherited from the prime zeta function, plus branch points
where the discriminant vanishes. Both parts deserve attention.

\subsection{What the tree zeta inherits}

The prime zeta function continues to the right half-plane by M\"obius inversion of
\eqref{eq:euler}:
\begin{equation}\label{eq:Pcont}
P(s) \;=\; \sum_{k \ge 1} \frac{\mu(k)}{k}\, \log \zeta(ks),
\qquad \operatorname{Re} s > 0,
\end{equation}
valid off the singular set. Equation \eqref{eq:Pcont} shows that $P$ has logarithmic
branch points at $s = 1/k$ (from the pole of $\zeta$) \emph{and at $s = \rho/k$ for
every nontrivial zero $\rho$ of $\zeta$}; these points cluster on the imaginary axis,
which is a natural boundary (Landau--Walfisz~\cite{landauwalfisz}). Two consequences:

\begin{enumerate}
\item Lifting arithmetic to trees does not dissolve the Riemann Hypothesis. The zeros
of $\zeta$ reappear, relocated, inside $Z_\cT$ as branch points at $\rho/k$.
\item No one-variable zeta of this family can contain more analytic information than
$P$ itself. Genuinely new invariants must either weight tree \emph{shape} (two-variable
zetas, Problem~\ref{prob:twovar}) or couple to the additive world (\S\ref{sec:bridge}).
\end{enumerate}

\subsection{Tree zeros}

The new singularities of $Z_\cT$ are the solutions of
\begin{equation}\label{eq:treezero}
P(s) = \tfrac14,
\end{equation}
square-root branch points of $Z_\cT$. We call them \emph{tree zeros} (they are the
zeros of the discriminant $1 - 4P$). On the real axis there is exactly one, $\sigma^{*}$,
since $P$ is strictly decreasing there. Off the axis:

\begin{lemma}\label{lem:strict}
For $\sigma > 0$ where $P(\sigma)$ converges and $t \neq 0$:
$|P(\sigma + it)| < P(\sigma)$.
\end{lemma}

\begin{proof}
$|\sum_p p^{-\sigma} p^{-it}| \le \sum_p p^{-\sigma}$ with equality iff all phases
$t \log p$ agree modulo $2\pi$. For $t \ne 0$ this fails already for $p = 2, 3$ since
$\log 3/\log 2$ is irrational.
\end{proof}

\begin{proposition}\label{prop:strip}
Every non-real tree zero satisfies $\operatorname{Re} s < \sigma^{*}$.
\end{proposition}

\begin{proof}
If $\sigma \ge \sigma^{*}$ and $t \neq 0$ then
$|P(\sigma+it)| < P(\sigma) \le P(\sigma^{*}) = \tfrac14$.
\end{proof}

The tree zeros with $0 < t < 50$ are listed in Table~\ref{tab:zeros} (method in
Appendix~\ref{app:verification}).

\begin{table}[ht]
\centering
\begin{tabular}{@{}cllc@{}}
\toprule
$k$ & $\operatorname{Re} s$ & $\operatorname{Im} s$ & $t \big/ \tfrac{2\pi}{\log 2}$ \\
\midrule
1 & $1.54427014392$ & $10.0122461654$ & $1.105$ \\
2 & $2.14220467429$ & $17.8224730193$ & $1.966$ \\
3 & $2.24736156780$ & $27.6431435747$ & $3.050$ \\
4 & $2.15502047591$ & $35.5407136013$ & $3.921$ \\
5 & $2.39048394209$ & $45.5394484588$ & $5.024$ \\
\bottomrule
\end{tabular}
\caption{All solutions of $P(s) = \tfrac14$ with $1.4 < \operatorname{Re} s < 2.7$,
$0 < \operatorname{Im} s < 50$ (each accompanied by its complex conjugate). A coarser
sweep of $1.05 < \operatorname{Re} s < 1.4$ found no further solutions in this height
range.}
\label{tab:zeros}
\end{table}

Two empirical observations. First, the ordinates hug integer multiples of
$2\pi/\log 2 = 9.0647\dots$: the term $2^{-s}$ dominates $P$, and $P(s) \approx \tfrac14$
requires $2^{-s}$ to return to the vicinity of its real value, which happens with that
period; the interference of $3^{-s}$ and later terms then shifts each zero (violently so
for $k=1$). Second, and decisively: \emph{the real parts wander}. The naive tree
analogue of the Riemann Hypothesis --- branch points on a vertical line --- is simply
false. The correct lifted question is distributional; see Problem~\ref{prob:zeros}.
By almost-periodicity of $P$ on vertical lines, the zero set is expected to have
positive density of ordinates and real parts dense in a subinterval of
$(\sigma_{\min}, \sigma^{*})$; identifying $\sigma_{\min}$ is open.

%======================================================================
\section{Thermodynamics of multiplicative microstates}\label{sec:thermo}

Write $E(T) = \log N(T)$. Then $N(T)^{-s} = e^{-sE(T)}$ and each corner's zeta becomes a
partition function at inverse temperature $\beta = s$:
\[
Z_{\mathcal X}(\beta) = \sum_{X \in \mathcal X} e^{-\beta E(X)}.
\]
For $\cM$ this is the primon gas of Julia~\cite{julia} and Spector~\cite{spector}: a
free bosonic gas with one mode per prime, energy $\log p$, partition function
$\zeta(\beta)$, and a Hagedorn catastrophe at $\beta = 1$, the setting made rigorous and
symmetry-rich by Bost--Connes~\cite{bostconnes}. The other three corners are gases of
\emph{structured} multiplicative states over the same spectrum.

\begin{center}
\begin{tabular}{@{}lllll@{}}
\toprule
gas & states & $\beta_c$ & $T_H = 1/\beta_c$ & transition \\
\midrule
multiset (primon) & prime multisets & $1$ & $1$ & pole \\
word & ordered factorizations & $1.399433\dots$ & $0.714575\dots$ & pole \\
bush & unordered trees & $1.988477\dots$ & $0.502897\dots$ & cusp, $Z_\cB \to 1$ \\
tree & planar trees & $2.597385\dots$ & $0.385003\dots$ & cusp, $Z_\cT \to \tfrac12$ \\
\bottomrule
\end{tabular}
\end{center}

The physics reading of Theorem~\ref{thm:ladder}: remembering order and grouping
multiplies the number of microstates at each energy (degeneracy
$g_{\mathcal X}(n)$), so the entropy grows faster and the gas can sustain less heat ---
the Hagedorn temperature drops from $1$ to $0.385$. The dichotomy of
Remark~\ref{rem:cusp} also becomes physical: for the associative gases the partition
function itself diverges at $\beta_c$; for the non-associative gases $Z$ stays finite
while the mean energy
\[
U(\beta) = -\frac{\partial}{\partial \beta} \log Z_\cT(\beta)
\;\sim\; \frac{\sqrt{|P'(\sigma^{*})|}}{2\,Z_\cT(\sigma^{*})}\,(\beta - \sigma^{*})^{-1/2}
\;\to\; \infty
\]
diverges --- a continuous (second-order-like) transition. In this precise sense,
\emph{associativity is a first-order axiom}.

%======================================================================
\section{The additive bridge: grouped quantities}\label{sec:bridge}

The additive world of the theory is flat: $\cA = \{\,n\,\ell : n \in \NN_{\ge1}\},$
heaps of a unit quantity $\ell$ (``one liter''), with
$+\colon \cA \times \cA \to \cA$ ordinary accumulation. The question left open by the
framework so far is how $\cT$ and $\cA$ interact. The answer proposed here: they do not
interact by a map $\cT \to \cA$. \emph{Multiplication acts on quantities}, and the
result of the action is neither a tree nor a flat quantity but a third kind of object.

\subsection{Grouped quantities}

\begin{definition}[Grouped quantities]\label{def:grouped}
Define \emph{packages} $t$ and \emph{quantities} $f$ by the mutual grammar
\[
t \;::=\; \ell \;\mid\; \pkg{f}\,,
\qquad\qquad
f \;::=\; t_1 \oplus t_2 \oplus \cdots \oplus t_m \quad (m \ge 1),
\]
where $\pkg{f}$ is a box containing the quantity $f$, and $\oplus$ is concatenation of
finite sequences of packages. Let $\cG$ denote the set of quantities. For $k \ge 2$
define the \emph{uniform grouping}
\[
k \odot f \;=\; \underbrace{\pkg{f} \oplus \pkg{f} \oplus \cdots \oplus \pkg{f}}_{k}.
\]
The \emph{flattening} $\alpha \colon \cG \to \cA$ counts atoms:
$\alpha(\ell) = \ell$, $\alpha(\pkg f) = \alpha(f)$,
$\alpha(f \oplus g) = \alpha(f) + \alpha(g)$. We identify $\cA$ with
$(\NN_{\ge1}, +)$ and write $\alpha(f) \in \NN_{\ge 1}$.
\end{definition}

Thus $\oplus$ is associative by construction (sequences), and $\cG$ refines $\cA$: a
quantity is a heap that remembers how it is packed. (One may further quotient $\cG$ by
reordering of packages; nothing below changes, as the proofs never use the order.)

\begin{definition}[The action]\label{def:action}
Define $\act \colon \cT \to \mathrm{Maps}(\cG, \cG)$ by
\[
\act(L_p)(f) = p \odot f,
\qquad
\act(T_1 \boxtimes T_2) = \act(T_1) \circ \act(T_2).
\]
\end{definition}

\begin{example}\label{ex:liters}
$\act(L_3)(\ell) = \pkg{\ell}\,\pkg{\ell}\,\pkg{\ell}$ --- three bottles of one liter.
Then
\[
\act(L_2 \boxtimes L_3)(\ell)
= 2 \odot \bigl(\pkg{\ell}\,\pkg{\ell}\,\pkg{\ell}\bigr)
= \pkg{\,\pkg{\ell}\,\pkg{\ell}\,\pkg{\ell}\,}\;\pkg{\,\pkg{\ell}\,\pkg{\ell}\,\pkg{\ell}\,}\,,
\]
two cartons each holding three bottles; whereas $\act(L_3 \boxtimes L_2)(\ell)$ is three
cartons each holding two bottles. Neither is ``six liters''; both flatten to
$6\ell$. The slogan of \S1 is now a statement about the image of $\act$: multiplication
outputs grouped quantities, never flat ones.
\end{example}

\subsection{The stratification theorems}

The three classical laws now separate into three different epistemic layers.

\begin{theorem}[Associativity is additively invisible --- exactly]\label{thm:assoc}
The action $\act$ factors through the associativity quotient $\cT \to \cW$, and the
induced map $\cW \to \mathrm{Maps}(\cG,\cG)$ is injective. Consequently the congruence
$\{(T, T') : \act(T) = \act(T')\}$ is precisely $\approx_a$.
\end{theorem}

\begin{proof}
\emph{Factoring:} $\mathrm{Maps}(\cG,\cG)$ is a monoid under composition, composition is
associative, and $\act$ sends $\boxtimes$ to $\circ$; hence
$\act((T_1 \boxtimes T_2)\boxtimes T_3) = \act(T_1)\circ\act(T_2)\circ\act(T_3)
= \act(T_1 \boxtimes(T_2 \boxtimes T_3))$, and $\act(T)$ depends only on
$\mathrm{word}(T)$.

\emph{Injectivity:} let $w = p_1 p_2 \cdots p_k$ and evaluate on the single atom:
$F = \act(w)(\ell) = p_1 \odot \act(p_2\cdots p_k)(\ell)$ consists of exactly $p_1$
identical top-level packages whose common content is $\act(p_2 \cdots p_k)(\ell)$.
So from $F$ one reads off $p_1$ (the number of top-level packages) and recurses on the
content; the recursion terminates at the boxless atom $\ell$. Distinct words therefore
give distinct maps.
\end{proof}

\begin{remark}[Why associativity feels inevitable]\label{rem:whyassoc}
Theorem~\ref{thm:assoc} explains a phenomenology: any semantics that interprets
multiplication as an \emph{iterable operation} --- anything of the form
$\boxtimes \mapsto \circ$ into an endomorphism monoid --- must identify at least
$\approx_a$, because composition of maps is associative whether we like it or not. Our
particular semantics identifies \emph{exactly} $\approx_a$: nothing more is lost.
Parenthesization is process history with no additive trace; to observe it, one would
have to make quantities out of operators (currying), which is precisely to leave the
additive world. By contrast, nothing in $\mathrm{Maps}(\cG,\cG)$ forces commutativity
--- and indeed it fails there.
\end{remark}

\begin{theorem}[The shadow of the action; commutativity dies only at the bottom]
\label{thm:shadowaction}
For all $T \in \cT$ and $f \in \cG$:
\[
\alpha\bigl(\act(T)(f)\bigr) \;=\; N(T)\cdot \alpha(f).
\]
Hence the induced action on flat quantities $\cA$ is multiplication by the shadow, and
two trees act identically on $\cA$ iff $N(T) = N(T')$, i.e.\ (by
Proposition~\ref{prop:fta}) iff $T \approx_{ac} T'$. In particular:
$\act(L_2 \boxtimes L_3) \neq \act(L_3 \boxtimes L_2)$ as maps of $\cG$, but they agree
after flattening.
\end{theorem}

\begin{proof}
Induction: $\alpha(p \odot f) = p\,\alpha(f)$ since boxes are $\alpha$-transparent, and
$\alpha(\act(T_1\boxtimes T_2)(f)) = N(T_1)\,\alpha(\act(T_2)(f)) =
N(T_1)N(T_2)\,\alpha(f)$. The characterization of equal flat actions is unique
factorization. Distinctness on $\cG$ is Theorem~\ref{thm:assoc} (different words).
\end{proof}

The two observation layers thus recover the two axes of the quotient square:
\emph{grouped} observation quotients by exactly associativity; \emph{flat} observation
quotients by exactly associativity and commutativity --- and the ``exactly'' of the
second statement is one more disguise of the Fundamental Theorem of Arithmetic.

\begin{theorem}[Distributivity is a flattening theorem]\label{thm:distrib}
For every $T \in \cT$ and $f, g \in \cG$:
\[
\act(T)(f \oplus g) \;\neq\; \act(T)(f) \oplus \act(T)(g)
\quad\text{in } \cG,
\qquad\text{yet}\qquad
\alpha\bigl(\act(T)(f \oplus g)\bigr) \;=\;
\alpha\bigl(\act(T)(f) \oplus \act(T)(g)\bigr).
\]
\end{theorem}

\begin{proof}
Let $p_1$ be the first letter of $\mathrm{word}(T)$. The left object consists of $p_1$
top-level packages; the right object consists of $2p_1$; since $p_1 \ge 2$ these
differ. Flattening both sides gives $N(T)(\alpha f + \alpha g)$ by
Theorem~\ref{thm:shadowaction}.
\end{proof}

``Multiply the merged heap'' and ``multiply separately, then merge'' are genuinely
different constructions with the same flat shadow --- the claim of the informal theory,
now a theorem. Note also the built-in one-sidedness: quantities act from the right of
nothing; there is no $(f \oplus g) \cdot T$ at all, because quantities are not
operators. It is striking that in Loday's arithmetree~\cite{loday}, where both arguments
\emph{are} trees, distributivity likewise survives on one side only.

\begin{center}
\begin{tabular}{@{}lccc@{}}
\toprule
law & in $\cT$ & acting on $\cG$ & after flattening $\alpha$ \\
\midrule
associativity & fails (free) & \textbf{holds} & holds \\
commutativity & fails (free) & fails & \textbf{holds} \\
distributivity & not formulable & fails & \textbf{holds} \\
\bottomrule
\end{tabular}
\end{center}

\subsection{Where the primes come from}\label{sec:primesfrom}

The theory so far imports its leaf alphabet from classical arithmetic. The bridge
removes this dependency: the leaves can be \emph{found} inside $\cG$.

\begin{definition}\label{def:groupable}
Call $n \in \NN_{\ge 2}$ \emph{groupable} if some quantity consisting of $k \ge 2$
identical packages, each containing $m \ge 2$ atoms, flattens to $n$ --- that is, if
$n = km$ with $k, m \ge 2$. Call $n$ \emph{prime} if it is not groupable.
\end{definition}

This is, of course, the classical definition of primality --- that is the point. It is
expressible entirely inside the grouped additive world, using only $\odot$ and
$\alpha$, with no imported number labels. One can therefore \emph{define} the leaf
alphabet as the set of non-groupable heap sizes and generate $\cT$ freely on it; the
numerical label of a leaf is recovered as the observable $\alpha$ of its heap. The
foundational question ``are prime leaves labeled by known primes, or can the labels be
derived?'' is thus answered: derived. What cannot be derived from the free multiplicative
structure alone (\S\ref{sec:priorart}) is derivable the moment the additive world and the
action are in place, because primality is an additive-observational property: a prime is
a heap size at which uniform regrouping is impossible.

%======================================================================
\section{Arithmetic functions on trees}\label{sec:functions}

Classical arithmetic functions pull back along $N$. For the Euler totient,
\[
\varphi_\cT(T) \;=\; N(T) \prod_{p \in \supp(T)} \Bigl(1 - \frac1p\Bigr)
\;=\; \varphi\bigl(N(T)\bigr)
\]
by Euler's product formula; e.g.\ $\varphi_\cT(L_2 \boxtimes L_3) = 6\cdot\tfrac12\cdot\tfrac23 = 2$.
But this pullback factors through the double quotient: it cannot distinguish
$L_2\boxtimes L_3$ from $L_3 \boxtimes L_2$. A genuinely tree-theoretic totient should
satisfy: (i) a recursion in $\boxtimes$; (ii) compatibility, i.e.\ it descends to
$\varphi \circ N$ after the double quotient; (iii) sensitivity, i.e.\ it separates some
pair of trees with equal shadow. No canonical candidate is known to us; we record this
as Problem~\ref{prob:totient}. The same design brief applies to M\"obius $\mu$,
divisor functions, and to two-variable refinements
$Z_\cT(s, q) = \sum_T q^{\,\mathrm{stat}(T)} N(T)^{-s}$ for shape statistics
$\mathrm{stat}$ (depth, right-branch count, Tamari rank): the interesting ones are
exactly those \emph{not} expressible through $P$ alone (\S\ref{sec:analytic}).

%======================================================================
\section{Open problems}\label{sec:open}

\begin{question}[Tree-zero distribution]\label{prob:zeros}
Describe the solution set of $P(s) = \tfrac14$ in $0 < \operatorname{Re} s < \sigma^{*}$:
density of ordinates up to height $T$; closure of the set of real parts; the left edge
$\sigma_{\min}$. Almost-periodicity of $P$ on vertical lines is the natural tool. More
generally, develop the ``spectral theory'' of the level sets $P(s) = c$ for
$0 < c < P(1^{+})$, of which the tree zeros are the case $c = \tfrac14$ and the word
poles the case $c = 1$.
\end{question}

\begin{question}[Tauberian rigor]\label{prob:tauberian}
Prove $\sum_{n \le x} g_\cT(n) \sim c\, x^{\sigma^{*}} (\log x)^{-3/2}$ (a Delange-type
theorem for Dirichlet series with a square-root singularity), and its analogue for
$g_\cB$.
\end{question}

\begin{question}[The cascade of the bush zeta]\label{prob:cascade}
The continuation $Z_\cB(s) = 1 - \sqrt{1 - 2P(s) - Z_\cB(2s)}$ propagates each
singularity of $Z_\cB$ at $s_0$ to potential singularities at $s_0/2, s_0/4, \dots$ and
mixes them with those of $P$. Map this self-similar singularity structure; determine
whether $Z_\cB$ has a natural boundary strictly to the right of
$\operatorname{Re} s = 0$.
\end{question}

\begin{question}[Beyond $P$]\label{prob:twovar}
Every one-variable zeta of \S\ref{sec:zeta} is a function of $P(s), P(2s), \dots$ Find a
natural two-variable zeta $Z_\cT(s,q)$ (weighting a shape statistic) that is provably
\emph{not} a function of the $P(ks)$, and study its singularities in $(s,q)$.
\end{question}

\begin{question}[Tree-sensitive totient]\label{prob:totient}
Construct $\varphi^{\mathrm{tree}}$ satisfying (i)--(iii) of \S\ref{sec:functions},
ideally with an inclusion--exclusion (M\"obius) interpretation over subtrees, perhaps
via the Connes--Kreimer-style coproduct on rooted trees.
\end{question}

\begin{question}[The full expression calculus]\label{prob:calculus}
\S\ref{sec:bridge} lets trees act on quantities but does not yet allow sums \emph{inside}
trees. Build the two-sorted expression grammar with both $\oplus$ and $\boxtimes$,
where distributivity becomes a directed rewrite rule (evaluation step) rather than an
equation, and relate its normal forms to Loday's grove addition~\cite{loday,brunoyasaki}.
\end{question}

\begin{question}[Axiomatics of the bridge]\label{prob:axioms}
Characterize the pair $(\cG, \alpha)$ abstractly: for which ``grouping worlds'' does
Theorem~\ref{thm:assoc} hold with \emph{exactly} $\approx_a$, and is $(\cG,\alpha)$
universal among them?
\end{question}

%======================================================================
\appendix
\section{Computational verification}\label{app:verification}

All computations are reproducible from the two scripts in the \texttt{verification/}
directory of this repository (\texttt{verify\_combinatorics.py}, standard library only;
\texttt{analytic.py}, requiring \texttt{mpmath} and \texttt{numpy}).

\subsection{Combinatorial checks}

For every $n \le 20000$: the closed formula of Proposition~\ref{prop:gT} agrees with a
direct recursive enumeration $g_\cT(n) = [n \in \mathbb{P}] + \sum_{d \mid n,\, 1<d<n}
g_\cT(d)\, g_\cT(n/d)$; the multinomial of Proposition~\ref{prop:gW} agrees with the
first-letter recursion; and the unordered recursion of Proposition~\ref{prop:gB}
reproduces the Wedderburn--Etherington numbers $1,1,1,2,3,6,11,23,46,98,207,451,983,2179$
at $n = 2^k$, $k \le 14$, as well as the hand-enumerated values $g_\cB(12)=2$,
$g_\cB(24)=4$, $g_\cB(30)=3$. The OEIS near-miss A375120 (which uses ordered pairs on
the diagonal) first diverges at $n = 144$: $42$ against the true $41$; a control query
confirmed the OEIS search interface finds A375120 from its own terms while the
$g_\cB$ terms return no match (July 2026).

\subsection{Analytic checks and constants}

With $g_\cT(n)$, $g_\cB(n)$ enumerated for $n \le 20000$, partial sums of the Dirichlet
series were compared with the closed/functional forms of Theorem~\ref{thm:funceq}:

\begin{center}\small
\begin{tabular}{@{}llll@{}}
\toprule
series & $s$ & partial sum vs.\ closed form & difference \\
\midrule
$Z_\cT$ & $3$ & $0.225547637865\ldots$ vs.\ $0.225705704214\ldots$ & $-1.6\cdot10^{-4}$ (tail) \\
$Z_\cT$ & $4$ & $0.084059063652\ldots$ vs.\ $0.084059066410\ldots$ & $-2.8\cdot10^{-9}$ \\
$Z_\cT$ & $6$ & agreement & $-3.1\cdot10^{-18}$ \\
$Z_\cB$ & $3$ & agreement & $-4.0\cdot10^{-7}$ \\
$Z_\cB$ & $4$ & agreement & $-1.1\cdot10^{-11}$ \\
\bottomrule
\end{tabular}
\end{center}

\noindent
the residual differences being exactly the truncation tails expected from the critical
exponents. Constants (via \texttt{mpmath.primezeta} at 25--30 significant digits):
\[
\begin{aligned}
\sigma^{*} &= 2.59738512713467162\dots, &\quad P'(\sigma^{*}) &= -0.230054805474\dots,\\
\sigma_\cB &= 1.98847685180090810\dots, &\quad Z_\cB(\sigma_\cB) &= 1\ (\text{verified to } 10^{-6}),\\
\sigma_\cW &= 1.39943332872633032\dots, &\quad P'(\sigma_\cW) &= -1.73115620186\dots,\\
\rho_{\text{Kalm\'ar}} &= 1.72864723899818362\dots, &\quad P(2) &= 0.45224742004106550\dots
\end{aligned}
\]

\subsection{Tree zeros}

The zeros of Table~\ref{tab:zeros} were located by a grid scan of
$|P_{\le 20000}(s) - \tfrac14|$ (truncated prime sum, step $0.02$) over
$[1.40, 2.70] \times [0.3, 50]$, followed by Newton polishing on the full
$\texttt{primezeta}$ at 20 significant digits; each reported zero satisfies
$|P(s) - \tfrac14| < 10^{-15}$. A coarser sweep of $[1.05, 1.40) \times [0.5, 50]$
(step $0.05 \times 0.25$) found no additional candidates; its only sub-threshold dip
flowed to zero \#1 of the table.

%======================================================================
\begin{thebibliography}{99}

\bibitem{bostconnes}
J.-B.~Bost and A.~Connes,
\emph{Hecke algebras, type III factors and phase transitions with spontaneous symmetry
breaking in number theory},
Selecta Math. (N.S.) \textbf{1} (1995), 411--457.

\bibitem{brunoyasaki}
A.~Bruno and D.~Yasaki,
\emph{The arithmetic of trees},
arXiv:0809.4448 (2008).

\bibitem{flajolet}
P.~Flajolet and R.~Sedgewick,
\emph{Analytic Combinatorics},
Cambridge University Press, 2009.

\bibitem{julia}
B.~L.~Julia,
\emph{Statistical theory of numbers},
in: Number Theory and Physics (Les Houches, 1989), Springer Proceedings in Physics
\textbf{47}, Springer, 1990, pp.~276--293.

\bibitem{kalmar}
L.~Kalm\'ar,
\emph{\"Uber die mittlere Anzahl der Produktdarstellungen der Zahlen},
Acta Litt. Sci. Szeged \textbf{5} (1931), 95--107.

\bibitem{landauwalfisz}
E.~Landau and A.~Walfisz,
\emph{\"Uber die Nichtfortsetzbarkeit einiger durch Dirichletsche Reihen definierter
Funktionen},
Rend. Circ. Mat. Palermo \textbf{44} (1920), 82--86.

\bibitem{loday}
J.-L.~Loday,
\emph{Arithmetree},
J. Algebra \textbf{267} (2003), 507--539; arXiv:math/0112034.

\bibitem{otter}
R.~Otter,
\emph{The number of trees},
Ann. of Math. \textbf{49} (1948), 583--599.

\bibitem{oeis}
N.~J.~A.~Sloane (ed.),
\emph{The On-Line Encyclopedia of Integer Sequences},
\url{https://oeis.org}. Entries A000108, A001190, A008480, A144757, A375120.

\bibitem{spector}
D.~Spector,
\emph{Supersymmetry and the M\"obius inversion function},
Comm. Math. Phys. \textbf{127} (1990), 239--252.

\bibitem{stanley}
R.~P.~Stanley,
\emph{Catalan Numbers},
Cambridge University Press, 2015.

\bibitem{tamari}
F.~M\"uller-Hoissen, J.~M.~Pallo, and J.~Stasheff (eds.),
\emph{Associahedra, Tamari Lattices and Related Structures},
Progress in Mathematics \textbf{299}, Birkh\"auser, 2012.

\end{thebibliography}

\end{document}
